\chapter{Evaluation}
\label{ch:evaluation}

This chapter specifies the evaluation protocol for the system. No full-study empirical result is reported in this thesis. The protocol defines the evidence required before comparing strategies.

\section{Methodology}

The current evaluation scope is a fixed stylized topology, segment-policy configuration, and attacker model. It cannot support claims about real lateral movement, topology- or vulnerability-density sensitivity, attacker-profile sensitivity, or general scalability. All current defensive actions have unit cost, so the budget compares equal action counts; it does not represent deployable, cost-aware segmentation. The current code supports the fixed baseline scenario and generic topology sources; the controlled topology-size tiers of the approved scale study are a research design and are not yet implemented as frozen study inputs.

A reproducible evaluation must provide a versioned runner; the exact topology, policy, vulnerability catalog, mission capabilities and required flows, entry point, and simulation configuration; the seed schedule; the selected action list, including random choices, and budget; complete per-trial outcomes; and the aggregation and statistical-analysis code. Each strategy must be run against the same declared scenario and equal action-count budget. The compared strategies are random selection, CVSS-priority ordering, topology-driven policy segmentation, simulation-informed patching, and simulated-annealing search over both action classes, each under equal action-count budgets of one through three with multiple fixed selection seeds.

The primary measure is the simulated mission impact across trials; blast radius is a secondary safety outcome computed from the same trials. Feasibility is compared on paired constrained and unconstrained plans, reporting pre-attack feasibility, unavailable required flows, and affected mission capabilities. Paired scenario and evaluation seeds are used so that compared plans differ only in the applied constraint. Runtime, memory use, and convergence may be reported only on a fixed documented environment after warm-up. If independently sampled groups are compared, the Mann--Whitney $U$ test~\cite{mann47} may be appropriate. Matched seeds or paired scenarios require a paired analysis selected before examining results.

\section{Results}

Results are pending the artifacts specified above. No percentage reduction, statistical conclusion, calibration claim, or comparison between strategies is supported until those artifacts are available.

\section{Strategy Comparison}

The protocol compares the available strategies on the declared fixed scenario, using the same entry point, seed schedule, and equal action-count budget. A comparison must report the full action list as well as mission-impact outcomes. It must not interpret lower modeled mission impact as evidence of real-world defensive effectiveness.

\section{Discussion}

The system is a context graph plus state-transition simulation informed by attack-graph literature, not a NetSPA or multiple-prerequisite attack-graph implementation. CVSS data describes severity characteristics. The simulator instead uses independently assigned, stylized success probabilities and uniformly selects eligible actions; it is neither calibrated to, nor a predictor of, real-world exploit likelihood.

The fixed stylized topology, policy, attacker model, mission capabilities, and unit action costs are deliberate model boundaries. They preclude conclusions about real attack propagation, cost-aware deployable segmentation, density or profile sensitivity, and scalability beyond explicitly measured scenarios. The contribution is a controlled simulation comparison of ranking strategies under a feasibility constraint; it is not a claim that multi-objective hardening or mission-impact assessment is new.
